
\def\PUDcodigo{EME117}

\def\PUDcodacad{04506.30}

\def\PUDdisciplina{Processos de Fabricação II}

\def\PUDsemestre{S6}

\def\PUDhoras{80}

%NOVOS ---------------------------------------------------
\def\PUDhorasTeoria{60}
\def\PUDhorasPratica{20}

\def\PUDinterECA{---}

\def\PUDatuaisECA{---}

\def\PUDrecursos{Projetor multimídia; Quadro branco e pincel; Aulas em laboratório.}

%----------------------------------------------------------

\def\PUDcreditos{4}

\def\PUDprerequisito{ \hyperref[PUD:EME108]{EME108} }

\def\PUDpreacad{ \hyperref[PUD:EME108]{04506.20} }

\def\PUDobjetivos{Conhecer os processos de usinagem voltados para a área de manutenção;  Compreender os parâmetros que influenciam nos processos de usinagem;  Conhecer as características das ferramentas de corte;  Saber avaliar a usinabilidade dos materiais;  Conhecer as características dos fluidos de corte.}

\def\PUDementa{Torneamento;  Fresamento;  Retificação;  Afiação;  Furação;  Movimentos e Grandezas nos Processos de Usinagem;  Mecanismos de Formação do Cavaco;  Forças e Potenciais de Corte;  Materiais para Ferramentas de Corte;  Avarias, Desgaste e Vida Útil das Ferramentas de Corte;  Usinabilidade dos Materiais;  Fluidos de Corte.}

\def\PUDconteudo{ & Equipamentos para tornear.  Processo de usinagem por torneamento. 
\\ 
 & Equipamentos para fresamento.  Processo de usinagem por fresamento. 
\\ 
 & Equipamentos para retificação.  Processo de usinagem por retificação. 
\\ 
 & Equipamentos para afiação.  Processo de usinagem por afiação. 
\\ 
 & Equipamentos para furação.  Processo de usinagem por furação. 
\\ 
 & Movimentos nos processos de usinagem.  Grandezas de avanço.  Grandezas de penetração.  Grandezas de corte. 
\\ 
 & Interface cavaco-ferramenta.  Controle da forma do cavaco.  Temperatura de corte. 
\\ 
 & Forças durante a usinagem.  Potências de usinagem.  Variações da força de corte.  Cálculo da pressão de corte.  Fatores que influenciam nas forças de avanço e de profundidade. 
\\ 
 & Classificação dos materiais para ferramentas.  Aços carbono para ferramentas.  Aços rápidos.  Ligas fundidas para ferramentas.  Metal duro.  Materiais cerâmicos.  Outros materiais para ferramentas. 
\\ 
 & Medição dos desgastes.  Mecanismos causadores do desgaste.  Fatores que influenciam no desgaste e vida útil da ferramenta.  Fatores de influência na rugosidade da peça. Curvas de vida útil da ferramenta.  Escolha dos parâmetros de usinagem. 
\\ 
 & Ensaios de usinabilidade.  Usinabilidade e propriedades do material.  Fatores metalúrgicos que afetam a usinabilidade. 
\\ 
 & Funções do fluido de corte.  Classificação dos fluidos de corte.  Seleção do fluido de corte.  Usinagem sem fluido de corte e com mínima quantidade de fluido. }

\def\PUDmetodo{Aulas expositórias que podem ser teóricas e/ou práticas, onde as práticas no laboratório serão marcadas no decorrer da disciplina.}

\def\PUDavalia{A avaliação será feita com aplicação de provas teóricas e/ou práticas, além da possibilidade de inclusão de trabalhos e seminários no decorrer da disciplina.}

\def\PUDbibbasica{ & \href{www.biblioteca.ifce.edu.br/index.asp?codigo_sophia=24286}{Ferraresi}, D.  Fundamentos da Usinagem dos Metais.  1ª ed.  São Paulo, SP: Edgard Blücher, 2009.  751p.  ISBN: 9788521202578. 
\\ 
 & \href{www.biblioteca.ifce.edu.br/index.asp?codigo_sophia=24076}{Diniz}, Anselmo Eduardo.  Tecnologia da usinagem dos materiais.  7ª ed.  São Paulo, SP: Artliber, 2010.  268p.  ISBN: 8587296019. 
\\ 
 & \href{www.biblioteca.ifce.edu.br/index.asp?codigo_sophia=23963}{Santos}, S. C.; Sales, W. F.  Aspectos Tribológicos da Usinagem dos Materiais.  1ª ed.  São Paulo, SP: Artliber, 2007.  246p.  ISBN: 9788588098381. }

\def\PUDbibcomplementar{ & \href{www.biblioteca.ifce.edu.br/index.asp?codigo_sophia=24061}{Chiaverini}, V.  Tecnologia Mecânica - Processos de Fabricação e Tratamentos.  Vol.2.  2ª ed.  São Paulo, SP: Pearson Makron Books, 1986.  315p.  ISBN: 9780074500903. 
\\ 
 &  \href{www.biblioteca.ifce.edu.br/index.asp?codigo_sophia=65142}{Machado}, A. R.; Abrão, A. M.; Coelho, R. T.; Silva, M. B.  Teoria da Usinagem dos Materiais.  2ª ed.  São Paulo, SP: Edgard Blücher, 2015.  407p.  ISBN: 9788521208464. 
\\ 
 &  \href{www.biblioteca.ifce.edu.br/index.asp?codigo_sophia=62257}{Fitzpatric}, M.  Introdução aos Processos de Usinagem.  1ª ed.  Porto Alegre, RS: AMGH, 2013.  488p.  ISBN: 9788580552287.
\\
 &  \href{www.biblioteca.ifce.edu.br/index.asp?codigo_sophia=67149}{Brasil}, Ministério da Educação.  Caderno de aulas práticas da tornearia.  Brasília, DF: Editora IFB, 2016.  103p.  ISBN: 9788564124424.
\\
 &  \href{www.biblioteca.ifce.edu.br/index.asp?codigo_sophia=63248}{Groover}, Mikell P.  Introdução aos processos de fabricação.  Rio de Janeiro, RJ: LTC, 2014.  737p.  ISBN: 9788521625193. }

\def\PUDautor{Francisco Nélio Costa Freitas}

\def\PUDdata{2013-05-23}

\def\PUDrevisao{1}

\def\PUDrevisor{Samuel Vieira Dias}

\def\PUDdatarevisao{2019-05-15}

\PUDcorpo
